\documentclass[conference,10pt]{IEEEtran}
\usepackage{amsmath,amssymb,booktabs,graphicx,algorithm,algpseudocode,multirow,float}
\usepackage{dblfloatfix}
\usepackage[section]{placeins}
\hyphenation{Leap-Frog}

\title{A Comparison of Stochastic Gradient Descent, Scaled Conjugate Gradient and LeapFrog Training of Feedforward Neural Networks}
\author{
\IEEEauthorblockN{Logan Shelver}
\IEEEauthorblockA{Student number: 26946777\\
RW741 Machine Learning\\
Department of Computer Science, Stellenbosch University\\
26946777@sun.ac.za}
\thanks{The code that produced the results reported here is available at https://github.com/naturalza/MLA-Assignment-3.}
}

\begin{document}
\maketitle

\begin{abstract}
Training a feedforward neural network is an unconstrained minimisation of an error defined on the network weights.
The study reported here compares stochastic gradient descent, scaled conjugate gradient, and the LeapFrog algorithm LFOP1(b) on three classification problems and three function-approximation problems of increasing complexity.
A single hidden layer is used.
The hidden-unit count and the control parameters of each trainer are selected by validation error.
A total of 30 independent runs and nonparametric statistical tests are used to rank the trainers.
Scaled conjugate gradient obtained the lowest mean test error on four of the six problems.
The three trainers were statistically indistinguishable on the Glass Identification problem.
\end{abstract}
\begin{IEEEkeywords}
neural networks, stochastic gradient descent, scaled conjugate gradient, LeapFrog, classification, function approximation
\end{IEEEkeywords}

\section{Introduction}
\label{sec:intro}
A feedforward neural network is trained by minimising a scalar error $E(\mathbf{w})$ over a packed weight vector $\mathbf{w}$.
The minimiser determines both the attained training error and the generalisation of the trained network.
First-order methods such as stochastic gradient descent (SGD) are simple to implement~\cite{lecun1998,rumelhart1986}.
The success of SGD depends on the learning rate, the momentum coefficient, and the mini-batch size~\cite{lecun1998,rumelhart1986}.
Second-order methods and dynamic methods reduce the dependence on the learning rate, the momentum coefficient, and the mini-batch size~\cite{moller1993,snyman1982}.

M{\o}ller introduced scaled conjugate gradient (SCG) as a conjugate-gradient trainer that avoids a line search by Levenberg--Marquardt scaling of a Hessian-vector product~\cite{moller1993}.
Snyman introduced LeapFrog as a dynamic unconstrained minimiser that treats $E$ as a potential energy~\cite{snyman1982}.
The leap-frog optimiser (LFOP) variant LFOP1(b) adds an adaptive time step~\cite{snyman1983}.
Engelbrecht discussed LeapFrog as a neural-network training algorithm~\cite{engelbrecht2007}.

The study reported here compares SGD, SCG, and LeapFrog LFOP1(b) as trainers of a one-hidden-layer feedforward network.
Three classification problems and three function-approximation problems of increasing complexity are used.
An empirical search determines the hidden-unit count of each problem.
A second empirical search determines the control parameters of each trainer.
A total of 30 independent initialisations and nonparametric tests are used to rank the trainers.

The remainder of the report is organised as follows.
Section~\ref{sec:background} describes the network model and the three training algorithms.
Section~\ref{sec:implementation} records the implementation choices.
Section~\ref{sec:empirical} defines the empirical process.
Section~\ref{sec:results} reports and discusses the results.
Section~\ref{sec:conclusions} concludes the report.

\section{Background}
\label{sec:background}
The three trainers compared in the study rest on different principles.
Section~\ref{sec:ffnn} describes the network.
Section~\ref{sec:sgd} describes stochastic gradient descent.
Section~\ref{sec:scg} describes scaled conjugate gradient.
Section~\ref{sec:leap} describes LeapFrog.
Section~\ref{sec:measures} then defines the performance measures and the statistical tests.

\subsection{Feedforward neural networks}
\label{sec:ffnn}
A one-hidden-layer network maps an input $\mathbf{x}\in\mathbb{R}^{n_{i}}$ to an output $\mathbf{y}\in\mathbb{R}^{n_{o}}$.
The hidden pre-activation and the hidden activation are
\begin{equation}
\mathbf{z}^{(1)} = \mathbf{W}^{(1)\top}\mathbf{x}+\mathbf{b}^{(1)}, \qquad
\mathbf{h} = \tanh(\mathbf{z}^{(1)}),
\end{equation}
where $\mathbf{W}^{(1)}\in\mathbb{R}^{n_{i}\times n_{h}}$ and $\mathbf{b}^{(1)}\in\mathbb{R}^{n_{h}}$.
The output pre-activation is
\begin{equation}
\mathbf{z}^{(2)} = \mathbf{W}^{(2)\top}\mathbf{h}+\mathbf{b}^{(2)}.
\end{equation}
Classification uses a softmax output.
Function approximation uses a linear output.
All weights and biases are packed into a single vector $\mathbf{w}\in\mathbb{R}^{N}$.
Training is then unconstrained minimisation of $E(\mathbf{w})$.
Bishop treated the packed-vector view of supervised learning as standard~\cite{bishop1995}.

The error $E$ is assumed continuously differentiable.
Reverse-mode backpropagation therefore yields $\nabla E(\mathbf{w})$.
Rumelhart \textit{et al.} derived the backpropagation recurrence~\cite{rumelhart1986}.

\subsection{Stochastic gradient descent}
\label{sec:sgd}
SGD replaces $\nabla E(\mathbf{w})$ by a gradient estimated on a mini-batch $B$ of the training set.
With momentum coefficient $\beta\in[0,1)$ and learning rate $\eta>0$, the update is
\begin{equation}
\mathbf{u}\leftarrow \beta\mathbf{u}-\eta\nabla E_{B}(\mathbf{w}), \qquad
\mathbf{w}\leftarrow\mathbf{w}+\mathbf{u}.
\end{equation}
Polyak introduced the momentum term as a heavy-ball correction~\cite{polyak1964}.
A batch size equal to the training-set size recovers full-batch gradient descent.
The learning rate, the momentum coefficient, and the batch size are control parameters~\cite{lecun1998}.
The control parameters are problem dependent.

\subsection{Scaled conjugate gradient}
\label{sec:scg}
Conjugate-gradient methods build search directions that are conjugate with respect to the Hessian of a quadratic model of $E$.
A line search is then required to choose a step length.
M{\o}ller avoided the line search by combining a conjugate-gradient direction with a Levenberg--Marquardt scalar $\lambda_{k}$ that keeps the model positive definite~\cite{moller1993}.

Let $\mathbf{r}_{k}=-\nabla E(\mathbf{w}_{k})$ and let $\mathbf{p}_{k}$ be the search direction.
A finite-difference Hessian-vector product
\begin{equation}
\mathbf{q}_{k}=\frac{\nabla E(\mathbf{w}_{k}+\sigma_{k}\mathbf{p}_{k})-\nabla E(\mathbf{w}_{k})}{\sigma_{k}},
\qquad \sigma_{k}=\frac{\sigma}{\lVert\mathbf{p}_{k}\rVert},
\end{equation}
is scaled to $\delta_{k}=\mathbf{p}_{k}^{\top}\mathbf{q}_{k}+(\lambda_{k}-\bar{\lambda}_{k})\lVert\mathbf{p}_{k}\rVert^{2}$.
If $\delta_{k}\le 0$, the scale $\lambda_{k}$ is raised until the model is positive definite.
The step length is $\alpha_{k}=\mu_{k}/\delta_{k}$ with $\mu_{k}=\mathbf{p}_{k}^{\top}\mathbf{r}_{k}$.
A comparison parameter $\Delta_{k}$ measures how well the quadratic model predicts the true reduction in $E$.
A successful step is accepted.
The scale $\lambda_{k}$ is lowered when $\Delta_{k}\ge 0.75$.
The scale $\lambda_{k}$ is raised when $\Delta_{k}<0.25$.
The search direction is then formed with the Polak--Ribi\`ere formula, or the search is restarted along $\mathbf{r}_{k+1}$.
M{\o}ller presented $\sigma\le 10^{-4}$ and $\lambda_{1}\le 10^{-6}$ as values that are not critical for success~\cite{moller1993}.
The present study still tunes $\sigma$ and $\lambda_{1}$ on a validation grid.

\subsection{LeapFrog unconstrained minimisation}
\label{sec:leap}
Snyman modelled unconstrained minimisation as the motion of a particle of unit mass in the potential $E$~\cite{snyman1982}.
The acceleration is $\mathbf{a}=-\nabla E(\mathbf{w})$.
A leap-frog integrator advances the position and the velocity,
\begin{equation}
\mathbf{w}_{k+1}=\mathbf{w}_{k}+\mathbf{v}_{k}\Delta t, \qquad
\mathbf{v}_{k+1}=\mathbf{v}_{k}+\mathbf{a}_{k+1}\Delta t.
\end{equation}
The displacement of a step is capped by a maximum length $\delta_{\max}$.
If the kinetic energy decreases, the particle is placed at the midpoint of the last step.
The velocity is then reduced, or set to zero after repeated decreases.
The LFOP1(b) modification grows $\Delta t$ after consecutive successful steps and halves $\Delta t$ after $m$ consecutive obtuse angles between successive accelerations~\cite{snyman1983}.
Snyman used $\Delta t=0.5$, $m=3$, and compared $\delta_{\max}=1$ with $\delta_{\max}=100$~\cite{snyman1983}.
The present study tunes $\Delta t$ and $\delta_{\max}$ on a validation grid and keeps $m=3$.

\subsection{Performance measures}
\label{sec:measures}
Classification is scored on a locked test set by accuracy, macro-averaged $F_{1}$, and mean categorical cross-entropy.
Function approximation is scored by mean squared error (MSE), root mean squared error, and the coefficient of determination $R^{2}$.
Training effort is recorded as the number of gradient evaluations and as wall-clock time.
The ranking of trainers uses test accuracy on classification problems and test MSE on function-approximation problems.

\subsection{Statistical comparison of stochastic trainers}
Random initial weights make each run a random variable.
Shapiro and Wilk tested normality of a sample~\cite{shapiro1965}.
If every sample of 30 test scores is consistent with normality, a one-way analysis of variance is used.
Kruskal and Wallis provided a nonparametric alternative when normality did not hold~\cite{kruskal1952}.
Pairwise differences are tested with the Mann--Whitney procedure~\cite{mann1947}.
Holm correction controls the family-wise error of the three pairwise tests on each problem~\cite{holm1979}.
Friedman ranking across problems is used as a further summary~\cite{friedman1937}.
A significance level of $0.05$ is used throughout.

SGD, SCG, and LeapFrog LFOP1(b) minimise the same differentiable error $E(\mathbf{w})$.
SGD is a first-order mini-batch method with three control parameters.
SCG is a scaled conjugate-gradient method with a finite-difference Hessian-vector product.
LeapFrog is a damped dynamic trajectory method with an adaptive time step.

\section{Implementation}
\label{sec:implementation}
The implementation records the architecture, the cost functions, the preprocessing, the three trainers, and the stopping rules.
Section~\ref{sec:arch} describes the network.
Section~\ref{sec:cost} describes the cost functions.
Section~\ref{sec:prep} describes the preprocessing.
Section~\ref{sec:sgd-impl} through Section~\ref{sec:stop} describe the trainers and the stopping rules.

\subsection{Network architecture and initialisation}
\label{sec:arch}
One hidden layer of hyperbolic-tangent units was used on every problem.
Classification used a softmax output.
Function approximation used a linear output.
Glorot and Bengio proposed the uniform Xavier distribution matched to hyperbolic-tangent units~\cite{glorot2010}.
Weights were drawn from the uniform Xavier distribution.
Biases were initialised at zero.
A packed vector $\mathbf{w}$ concatenates $\mathbf{W}^{(1)}$, $\mathbf{b}^{(1)}$, $\mathbf{W}^{(2)}$, and $\mathbf{b}^{(2)}$ in that order.

\subsection{Cost functions}
\label{sec:cost}
Classification used mean categorical cross-entropy
\begin{equation}
E(\mathbf{w})=-\frac{1}{n}\sum_{i=1}^{n}\sum_{c=1}^{n_{o}} y_{ic}\log(\hat{y}_{ic}+\varepsilon),
\end{equation}
with $\varepsilon=10^{-12}$.
Function approximation used MSE
\begin{equation}
E(\mathbf{w})=\frac{1}{n\,n_{o}}\sum_{i=1}^{n}\lVert\hat{\mathbf{y}}_{i}-\mathbf{y}_{i}\rVert^{2}.
\end{equation}
The analytic gradient matched $E$.
The gradient was checked against a central finite difference on exclusive-or.
The relative error of the finite-difference check was $5.4\times 10^{-11}$.

\subsection{Data preprocessing}
\label{sec:prep}
Input features were standardised from the training split only.
The training mean and the training standard deviation were applied to the validation split and to the test split.
Classification targets were one-hot encoded.
A linear output with MSE is sensitive to target scale.
Regression targets were therefore standardised from the training split.
Splits were stratified on classification problems.
A random split was used on function-approximation problems.
The outer proportions were $70\%$ train, $15\%$ validation, and $15\%$ test, with a fixed seed.
Every trainer therefore saw the same partitions.

\subsection{Stochastic gradient descent implementation}
\label{sec:sgd-impl}
Algorithm~\ref{alg:sgd} states the mini-batch update.
A full pass over the training set is one epoch.
The validation error was recorded at the end of each epoch.

\begin{algorithm}[t]
\caption{Mini-batch stochastic gradient descent}
\label{alg:sgd}
\begin{algorithmic}[1]
\Require packed weights $\mathbf{w}$, learning rate $\eta$, momentum $\beta$, batch size $s$
\State set the momentum buffer $\mathbf{u}$ to zero
\For{every epoch}
  \State shuffle the training set
  \For{every mini-batch $B$ of size $s$}
    \State set $\mathbf{u}$ to $\beta\mathbf{u}-\eta\nabla E_{B}(\mathbf{w})$
    \State set $\mathbf{w}$ to $\mathbf{w}+\mathbf{u}$
  \EndFor
  \State record the training error and the validation error
\EndFor
\end{algorithmic}
\end{algorithm}

\subsection{Scaled conjugate gradient implementation}
Algorithm~\ref{alg:scg} follows the numbered steps of M{\o}ller~\cite{moller1993}.
Two implementation safeguards were added.
A search direction that is not a descent direction was replaced by the steepest-descent vector.
A failed reduction restarted conjugacy from steepest descent after $\lambda_{k}$ was raised.
The restart prevented a near-singular conjugate direction from stalling the iteration.
The two safeguards did not change the scale-update formulae of M{\o}ller.

\begin{algorithm}[t]
\caption{Scaled conjugate gradient}
\label{alg:scg}
\begin{algorithmic}[1]
\Require packed weights $\mathbf{w}$, interval $\sigma$, scale $\lambda$
\State set $\mathbf{r}$ to $-\nabla E(\mathbf{w})$, set $\mathbf{p}$ to $\mathbf{r}$, set $\bar{\lambda}$ to $0$, and mark the iteration as successful
\While{the stopping rule is not met}
  \If{the previous iteration succeeded}
    \State estimate $\mathbf{q}$ and set $\delta_{k}$ to $\mathbf{p}^{\top}\mathbf{q}$
  \EndIf
  \State set $\delta_{k}$ to $\delta_{k}+(\lambda-\bar{\lambda})\lVert\mathbf{p}\rVert^{2}$
  \If{$\delta_{k}\le 0$}
    \State raise $\lambda$ until the quadratic model is positive definite
  \EndIf
  \State set $\alpha$ to $(\mathbf{p}^{\top}\mathbf{r})/\delta_{k}$
  \State compute $\Delta_{k}$ from the true reduction and the quadratic reduction
  \If{$\Delta_{k}\ge 0$}
    \State set $\mathbf{w}$ to $\mathbf{w}+\alpha\mathbf{p}$, update $\mathbf{r}$, and accept the step
    \State restart the direction or apply the Polak--Ribi\`ere formula for $\mathbf{p}$
  \Else
    \State reject the step, raise $\lambda$, and set $\mathbf{p}$ to $\mathbf{r}$
  \EndIf
\EndWhile
\end{algorithmic}
\end{algorithm}

\subsection{LeapFrog implementation}
Algorithm~\ref{alg:leap} follows algorithm (23) of Snyman~\cite{snyman1982} with the LFOP1(b) time-step rules~\cite{snyman1983}.
The time step was clipped to $[10^{-8},10]$ to prevent overflow after many consecutive successful steps.

\begin{algorithm}[t]
\caption{LeapFrog LFOP1(b)}
\label{alg:leap}
\begin{algorithmic}[1]
\Require packed weights $\mathbf{w}$, time step $\Delta t$, maximum step $\delta_{\max}$, and $m=3$
\State set $\mathbf{a}$ to $-\nabla E(\mathbf{w})$ and set $\mathbf{v}$ to $\tfrac{1}{2}\mathbf{a}\Delta t$
\While{the stopping rule is not met}
  \State cap $\lVert\mathbf{v}\rVert\Delta t$ by $\delta_{\max}$
  \If{the last $m$ successive accelerations form obtuse angles}
    \State halve $\Delta t$ and restart from the midpoint with reduced velocity
  \EndIf
  \State set $\mathbf{w}$ to $\mathbf{w}+\mathbf{v}\Delta t$, set $\mathbf{a}$ to $-\nabla E(\mathbf{w})$, and set $\mathbf{v}$ to $\mathbf{v}+\mathbf{a}\Delta t$
  \If{kinetic energy decreased}
    \State move to the midpoint and reduce or zero the velocity
  \EndIf
\EndWhile
\end{algorithmic}
\end{algorithm}

\subsection{Stopping criteria}
\label{sec:stop}
Every trainer stopped at the first of three events: a validation-error plateau of 60 recorded steps, a gradient norm below $10^{-8}$, or a cap of $20\,000$ gradient evaluations.
The weights with the lowest validation error were restored.
The same three events were used for all algorithms.
A comparison of test error was therefore not confounded by unequal stopping rules.

A shared network and a shared analytic gradient were used by SGD, SCG, and LeapFrog.
The three algorithms differed only in the update of $\mathbf{w}$.
Python~3 and NumPy were used.
No third-party implementation of SCG or LeapFrog was used for the reported numbers.

\section{Empirical Process}
\label{sec:empirical}
The empirical process covers the problems, the hidden-unit search, the control-parameter search, and the 30-run comparison.
Section~\ref{sec:problems} describes the problems.
Section~\ref{sec:hu} describes the hidden-unit search.
Section~\ref{sec:cp} describes the control-parameter search.
Section~\ref{sec:final} describes the final comparison.

\subsection{Problems}
\label{sec:problems}
Three classification problems and three function-approximation problems were selected to form a complexity ladder.
Table~\ref{tab:problems} records the dimensions after the split.

\begin{table}[H]
\caption{Problem dimensions after the $70/15/15$ split.}
\label{tab:problems}
\begin{center}
\begin{tabular}{lrrrrr}
\toprule
Problem & $n_{i}$ & $n_{o}$ & Train & Val. & Test \\
\midrule
Iris & 4 & 3 & 104 & 23 & 23 \\
Breast Cancer & 30 & 2 & 397 & 86 & 86 \\
Glass & 9 & 6 & 149 & 32 & 33 \\
sinc & 1 & 1 & 139 & 31 & 30 \\
sincos & 2 & 1 & 279 & 61 & 60 \\
Friedman \#1 & 10 & 1 & 279 & 61 & 60 \\
\bottomrule
\end{tabular}
\end{center}
\end{table}

The Iris data set has 150 samples, four continuous features, and three flower classes.
Iris was treated as the easy classification problem.
The Wisconsin Breast Cancer Diagnostic data set has 569 samples, 30 real-valued features, and two classes.
Breast Cancer was treated as the medium classification problem.
The Glass Identification data set has 214 samples, nine chemical features, and six imbalanced classes.
Glass was treated as the hard classification problem.
The three classification sets were taken from the University of California, Irvine (UCI) repository.\footnote{https://archive.ics.uci.edu}

A one-dimensional sinc function $f$ was treated as the easy function-approximation problem.
The function $f$ satisfies $f(x)=\sin(x)/x$ on $[-10,10]$.
A total of 200 noise-free samples were drawn.
A two-dimensional function $g$ was treated as the medium function-approximation problem.
The function $g$ satisfies $g(x,y)=\sin(x)\cos(y)$ on $[-\pi,\pi]^{2}$.
A total of 400 noise-free samples were drawn.
Friedman \#1 was treated as the hard function-approximation problem~\cite{friedman1991}.
The problem has ten inputs.
Five inputs are informative.
Five inputs are irrelevant.
A total of 400 samples were drawn with Gaussian noise of standard deviation $0.1$.

\subsection{Hidden-unit search}
\label{sec:hu}
Optimality of the hidden-unit count was defined as follows.
Ten independent runs were performed at each candidate count.
The selected count was the smallest count with a mean validation error within five percent of the best mean on the problem.
The five-percent band preferred a smaller network when the validation error was statistically close to the best value.
SCG has few control parameters.
SCG converges in few iterations.
SCG with $\sigma=10^{-4}$ and $\lambda_{1}=10^{-6}$ was therefore used for the architecture search.
The grid $\{2,4,6,8,12,16\}$ was used on Iris and sinc.
The grid $\{4,8,12,16,24\}$ was used on sincos.
The grid $\{4,8,12,16,24,32\}$ was used on Breast Cancer, Glass, and Friedman \#1.
The selected count was then frozen for every trainer on that problem.

\subsection{Control-parameter search}
\label{sec:cp}
After the architecture was locked, each trainer was tuned on each problem by mean validation error over eight runs.
The SGD grid was $\eta\in\{0.001,0.01,0.1\}$, $\beta\in\{0.0,0.9\}$, and batch size in $\{16,32,n\}$.
The SCG grid was $\sigma\in\{10^{-4},10^{-5}\}$ and $\lambda_{1}\in\{10^{-6},10^{-4}\}$.
The LeapFrog grid was $\Delta t\in\{0.1,0.5,1.0\}$ and $\delta_{\max}\in\{1,10,100\}$.
The value of $m$ was held at three.
The winning tuple of each trainer on each problem was locked for the final comparison.

\subsection{Final comparison}
\label{sec:final}
A total of 30 independent Xavier initialisations were drawn for each pair of algorithm and problem.
The splits, the hidden-unit counts, and the control parameters were held fixed.
Test accuracy was the primary classification score.
Test MSE was the primary function-approximation score.
Shapiro--Wilk tests were applied to the 30 primary scores.
Kruskal--Wallis tests and Holm-corrected Mann--Whitney tests were then applied on each problem.
A Friedman test ranked the three trainers across the six problems.

The protocol had four locked stages: problem construction, hidden-unit search, control-parameter search, and a 30-run comparison.
Reported results were produced only from the locked protocol.

\section{Results and Discussion}
\label{sec:results}
The results address the hidden-unit search, the selection of the control parameters, and the comparison of the three trainers.
Section~\ref{sec:res-hu} reports the hidden-unit search.
Section~\ref{sec:res-cp} reports the tuning.
Section~\ref{sec:res-clf} and Section~\ref{sec:res-reg} compare the trainers on classification and on function approximation.
Section~\ref{sec:res-stat} reports the statistical tests.

\subsection{Hidden-unit search}
\label{sec:res-hu}
Table~\ref{tab:hidden} records the selected hidden-unit counts.

\begin{table}[H]
\caption{Selected hidden-unit counts from the validation search.}
\label{tab:hidden}
\begin{center}
\begin{tabular}{lrrr}
\toprule
Problem & Hidden units & Val.\ error & Train error \\
\midrule
Iris & 4 & 0.0316 & 0.0340 \\
Breast Cancer & 32 & 0.0024 & 0.0000 \\
Glass & 4 & 0.7921 & 0.4507 \\
sinc & 8 & 0.0043 & 0.0036 \\
sincos & 24 & 0.0184 & 0.0114 \\
Friedman \#1 & 12 & 0.0262 & 0.0055 \\
\bottomrule
\end{tabular}
\end{center}
\end{table}

Iris selected four hidden units.
Figure~\ref{fig:hidden-iris} shows the validation curve.
The mean validation cross-entropy at four units was $0.0316$.
Larger nets did not improve on that value by more than five percent.
\begin{figure}[H]
\centering
\includegraphics[width=\columnwidth]{figures/hidden_iris.pdf}
\caption{Hidden-unit search on Iris.}
\label{fig:hidden-iris}
\end{figure}
Breast Cancer selected 32 hidden units.
Figure~\ref{fig:hidden-breast} shows the validation curve.
Training cross-entropy reached zero from eight units upward.
Validation cross-entropy continued to fall.
The five-percent rule therefore retained the largest grid value.
\begin{figure}[H]
\centering
\includegraphics[width=\columnwidth]{figures/hidden_breast_cancer.pdf}
\caption{Hidden-unit search on Breast Cancer.}
\label{fig:hidden-breast}
\end{figure}
Glass selected four hidden units.
Figure~\ref{fig:hidden-glass} shows the validation curve.
The lowest mean validation error on Glass was $0.776$ at eight units.
A count of four units lay inside the five-percent band and was therefore preferred.
Training error on Glass kept falling while the hidden layer grew.
Validation error rose.
The pattern is the signature of overfitting on a small, imbalanced set.
\begin{figure}[H]
\centering
\includegraphics[width=\columnwidth]{figures/hidden_glass.pdf}
\caption{Hidden-unit search on Glass.}
\label{fig:hidden-glass}
\end{figure}

The sinc problem selected eight hidden units.
The mean validation MSE was $0.0043$.
Figure~\ref{fig:hidden-sinc} shows the validation curve.
\begin{figure}[H]
\centering
\includegraphics[width=\columnwidth]{figures/hidden_sinc.pdf}
\caption{Hidden-unit search on sinc.}
\label{fig:hidden-sinc}
\end{figure}
The sincos problem selected 24 hidden units.
Figure~\ref{fig:hidden-sincos} shows the validation curve.
\begin{figure}[H]
\centering
\includegraphics[width=\columnwidth]{figures/hidden_sincos.pdf}
\caption{Hidden-unit search on sincos.}
\label{fig:hidden-sincos}
\end{figure}
Friedman \#1 selected 12 hidden units.
Figure~\ref{fig:hidden-friedman} shows the validation curve.
The mean validation MSE at 12 units was $0.0262$.
Both smaller nets and larger nets were worse.
\begin{figure}[H]
\centering
\includegraphics[width=\columnwidth]{figures/hidden_friedman1.pdf}
\caption{Hidden-unit search on Friedman \#1.}
\label{fig:hidden-friedman}
\end{figure}


\subsection{Control-parameter search}
\label{sec:res-cp}
Table~\ref{tab:params} records the winning control parameters and the mean validation error of each winner.

\begin{table}[H]
\caption{Locked control parameters and mean validation error.}
\label{tab:params}
\begin{center}
\begin{tabular}{llr}
\toprule
Problem and trainer & Locked parameters & Val.\ error \\
\midrule
Iris, SGD & $\eta=0.1$, $\beta=0.9$, batch $32$ & 0.031 \\
Iris, SCG & $\sigma=10^{-5}$, $\lambda_{1}=10^{-4}$ & 0.023 \\
Iris, LeapFrog & $\Delta t=1.0$, $\delta_{\max}=1$ & 0.070 \\
Breast Cancer, SGD & $\eta=0.1$, $\beta=0.9$, batch $32$ & 0.008 \\
Breast Cancer, SCG & $\sigma=10^{-4}$, $\lambda_{1}=10^{-6}$ & 0.001 \\
Breast Cancer, LeapFrog & $\Delta t=1.0$, $\delta_{\max}=10$ & 0.006 \\
Glass, SGD & $\eta=0.1$, $\beta=0.9$, batch $16$ & 0.692 \\
Glass, SCG & $\sigma=10^{-4}$, $\lambda_{1}=10^{-6}$ & 0.846 \\
Glass, LeapFrog & $\Delta t=1.0$, $\delta_{\max}=10$ & 0.830 \\
sinc, SGD & $\eta=0.1$, $\beta=0.9$, batch $32$ & 0.022 \\
sinc, SCG & $\sigma=10^{-4}$, $\lambda_{1}=10^{-4}$ & 0.004 \\
sinc, LeapFrog & $\Delta t=1.0$, $\delta_{\max}=1$ & 0.010 \\
sincos, SGD & $\eta=0.01$, $\beta=0.9$, batch $16$ & 0.041 \\
sincos, SCG & $\sigma=10^{-5}$, $\lambda_{1}=10^{-4}$ & 0.017 \\
sincos, LeapFrog & $\Delta t=1.0$, $\delta_{\max}=1$ & 0.035 \\
Friedman \#1, SGD & $\eta=0.1$, $\beta=0.0$, batch $16$ & 0.074 \\
Friedman \#1, SCG & $\sigma=10^{-5}$, $\lambda_{1}=10^{-6}$ & 0.017 \\
Friedman \#1, LeapFrog & $\Delta t=1.0$, $\delta_{\max}=1$ & 0.014 \\
\bottomrule
\end{tabular}
\end{center}
\end{table}

SGD preferred $\eta=0.1$ and $\beta=0.9$ on four of the six problems.
Full-batch SGD was never selected.
SCG preferred values at the defaults of M{\o}ller or at the neighbouring grid point on every problem.
The observation agrees with the claim that $\sigma$ and $\lambda_{1}$ are not critical~\cite{moller1993}.
LeapFrog preferred $\Delta t=1.0$ on every problem.
The maximum step $\delta_{\max}$ was $1$ on Iris, sinc, sincos, and Friedman \#1.
The maximum step $\delta_{\max}$ was $10$ on Breast Cancer and Glass.
A value $\delta_{\max}=100$ with $\Delta t=1.0$ produced a mean validation error of $0.67$ on sincos.
An oversized step is therefore harmful.

On Glass, SGD obtained a mean validation cross-entropy of $0.692$.
SCG obtained $0.846$ on Glass.
LeapFrog obtained $0.830$ on Glass.
The architecture of four hidden units, chosen by SCG, is a small model of a six-class problem.
The lower validation error of SGD on Glass is consistent with mini-batch regularisation on a small architecture.


\subsection{Classification test performance}
\label{sec:res-clf}
Figure~\ref{fig:final-clf} and Table~\ref{tab:clf} report mean test accuracy over 30 runs.

\begin{figure}[H]
\centering
\includegraphics[width=\columnwidth]{figures/final_classification.pdf}
\caption{Mean test accuracy on the classification problems.}
\label{fig:final-clf}
\end{figure}

\begin{table}[H]
\caption{Mean test accuracy $\pm$ standard deviation over 30 runs.}
\label{tab:clf}
\begin{center}
\begin{tabular}{lccc}
\toprule
Problem & SGD & SCG & LeapFrog \\
\midrule
Iris & $0.951\pm 0.027$ & $0.988\pm 0.022$ & $0.965\pm 0.041$ \\
Breast Cancer & $0.975\pm 0.004$ & $0.973\pm 0.009$ & $0.971\pm 0.007$ \\
Glass & $0.617\pm 0.071$ & $0.619\pm 0.072$ & $0.615\pm 0.083$ \\
\bottomrule
\end{tabular}
\end{center}
\end{table}

SCG obtained the highest mean accuracy on Iris, equal to $0.988$.
SGD obtained the highest mean accuracy on Breast Cancer, equal to $0.975$.
The three means on Breast Cancer lay inside one standard deviation of each other.
Glass remained difficult for every trainer.
Mean test accuracy was $0.617$ for SGD, $0.619$ for SCG, and $0.615$ for LeapFrog.
The standard deviations were $0.071$, $0.072$, and $0.083$.
Macro-averaged $F_{1}$ on Glass separated the trainers more clearly than accuracy.
The macro-averaged $F_{1}$ was $0.499$ for SCG, $0.484$ for SGD, and $0.435$ for LeapFrog.
The small hidden layer, the six-class imbalance, and the 33-sample test set explained the high variance on Glass.


\subsection{Function-approximation test performance}
\label{sec:res-reg}
Figure~\ref{fig:final-reg} and Table~\ref{tab:reg} report mean test MSE over 30 runs.

\begin{figure}[H]
\centering
\includegraphics[width=\columnwidth]{figures/final_regression.pdf}
\caption{Mean test mean squared error on the function-approximation problems.}
\label{fig:final-reg}
\end{figure}

\begin{table}[H]
\caption{Mean test mean squared error $\pm$ standard deviation over 30 runs.}
\label{tab:reg}
\begin{center}
\begin{tabular}{lccc}
\toprule
Problem & SGD & SCG & LeapFrog \\
\midrule
sinc & $0.022\pm 0.045$ & $0.001\pm 0.004$ & $0.005\pm 0.008$ \\
sincos & $0.021\pm 0.007$ & $0.005\pm 0.006$ & $0.014\pm 0.008$ \\
Friedman \#1 & $0.044\pm 0.045$ & $0.014\pm 0.052$ & $0.062\pm 0.099$ \\
\bottomrule
\end{tabular}
\end{center}
\end{table}

SCG obtained the lowest mean test MSE on sinc at $0.0012$, on sincos at $0.0051$, and on Friedman \#1 at $0.0141$.
LeapFrog was second on sinc and sincos.
SGD had a large standard deviation on sinc of $0.045$ relative to SCG at $0.004$.
The gap is consistent with residual sensitivity of SGD to random initialisation.
LeapFrog was the weakest of the three methods on Friedman \#1, with mean MSE $0.062$ and standard deviation $0.099$.
Mean coefficients of determination on sinc were $0.998$ for SCG, $0.993$ for LeapFrog, and $0.966$ for SGD.
Irrelevant inputs and additive noise produced a less smooth error surface.
The kinetic-energy heuristic of LeapFrog was less reliable on the Friedman \#1 error surface than the scaled curvature of SCG.

\subsection{Statistical comparison}
\label{sec:res-stat}
Table~\ref{tab:stats} records the omnibus test and the Holm-corrected pairwise decisions.

\begin{table}[H]
\caption{Omnibus tests on the 30 primary test scores.}
\label{tab:stats}
\begin{center}
\begin{tabular}{llcc}
\toprule
Problem & Test & $p$ & Decision at $0.05$ \\
\midrule
Iris & Kruskal--Wallis & $3.5\times 10^{-5}$ & reject \\
Breast Cancer & Kruskal--Wallis & $0.073$ & do not reject \\
Glass & analysis of variance & $0.98$ & do not reject \\
sinc & Kruskal--Wallis & $1.5\times 10^{-6}$ & reject \\
sincos & Kruskal--Wallis & $4.5\times 10^{-11}$ & reject \\
Friedman \#1 & Kruskal--Wallis & $2.5\times 10^{-10}$ & reject \\
All six means & Friedman & $0.042$ & reject \\
\bottomrule
\end{tabular}
\end{center}
\end{table}

Shapiro--Wilk tests rejected normality on five of the six problems.
Kruskal--Wallis tests were therefore used on those problems.
Glass was consistent with normality.
A one-way analysis of variance was therefore used on Glass.

The Kruskal--Wallis test rejected equality of the three trainers on Iris, with $p=3.5\times 10^{-5}$.
The same test rejected equality on sinc, with $p=1.5\times 10^{-6}$.
The same test rejected equality on sincos, with $p=4.5\times 10^{-11}$.
The same test rejected equality on Friedman \#1, with $p=2.5\times 10^{-10}$.
Holm-corrected Mann--Whitney tests found SCG better than SGD and better than LeapFrog on Iris, sinc, sincos, and Friedman \#1.
SGD was better than LeapFrog on sinc, sincos, and Friedman \#1.
Breast Cancer did not reject equality, with a Kruskal--Wallis $p$-value of $0.073$.
Glass did not reject equality, with an analysis-of-variance $p$-value of $0.98$.

A Friedman test on the six mean primary scores rejected equality of ranks, with $p=0.042$.
The average ranks were $1.17$ for SCG, $2.33$ for SGD, and $2.50$ for LeapFrog, with rank one denoting the best mean.

The ranking was therefore not uniform across problem types.
On smooth, low-dimensional function approximation, SCG was the strongest of the three methods.
On linearly well-separated medical data with 30 features, mini-batch SGD matched SCG.
On a small imbalanced multi-class problem with a parsimonious hidden layer, no trainer dominated.

\subsection{Discussion}
Three observations follow from the tables.
First, the claim of M{\o}ller that SCG has no critical user parameters is supported by the validation grid~\cite{moller1993}.
Neighbouring $(\sigma,\lambda_{1})$ pairs produced similar validation errors.
Second, SGD remained competitive when $\eta$ and $\beta$ were tuned per problem.
SGD used far more gradient evaluations than SCG and LeapFrog on every problem.
Each mini-batch step is one evaluation.
Early stopping is checked only once per epoch.
Third, LeapFrog LFOP1(b) was an effective neural-network trainer on the easy and medium problems.
Engelbrecht discussed LeapFrog for neural-network training~\cite{engelbrecht2007}.
LeapFrog was more sensitive to $\delta_{\max}$ than SCG was to $\sigma$.
LeapFrog was the least stable of the three methods on Friedman \#1.

A limitation of the study was the use of a single hidden layer and of small data sets.
A second limitation was the architecture search with SCG only.
A hidden-unit count that is optimal for SCG need not be optimal for SGD or LeapFrog.
The protocol matched the stated requirement of one hidden-layer size per problem, not per algorithm.
SCG was the most reliable of the three trainers across the six problems.
SGD matched SCG on Breast Cancer after a per-problem grid search.
LeapFrog was competitive on classification and on the two smoother regression problems.
LeapFrog was the weakest method on Friedman \#1.
Glass defeated all three trainers to a similar degree.


\section{Conclusions}
\label{sec:conclusions}
Stochastic gradient descent (SGD), scaled conjugate gradient (SCG), and the LeapFrog algorithm LFOP1(b) were compared as trainers of a one-hidden-layer feedforward neural network.
Six problems of increasing complexity were used.
Hidden-unit counts and control parameters were selected by validation error.
A total of 30 independent runs supplied the test scores.

SCG obtained the lowest mean test error on Iris, sinc, sincos, and Friedman \#1.
SGD obtained the lowest mean test accuracy on Breast Cancer.
All three means on Breast Cancer were close.
The three trainers were indistinguishable on Glass Identification at the reported precision.
The study therefore supported SCG as the default trainer of the three methods.
A tuned SGD remained adequate on well-separated classification data.
LeapFrog required a careful cap on the step length.

The comparison was restricted to one hidden layer, to six small problems, and to a single architecture per problem.
The listed restrictions bounded the claim to the protocol that was run.

\bibliographystyle{IEEEtranS}
\bibliography{references}

\appendix
\section{Acronyms}
\begin{description}
\item[LFOP] leap-frog optimiser
\item[MSE] mean squared error
\item[SCG] scaled conjugate gradient
\item[SGD] stochastic gradient descent
\item[UCI] University of California, Irvine
\end{description}

\section{Symbols}
\begin{description}
\item[$\mathbf{a}$] LeapFrog acceleration, $\mathbf{a}=-\nabla E(\mathbf{w})$
\item[$\alpha_{k}$] SCG step length
\item[$\mathbf{b}$] bias vector
\item[$\beta$] SGD momentum coefficient
\item[$B$] mini-batch
\item[$\delta_{k}$] SCG scaled curvature
\item[$\delta_{\max}$] LeapFrog maximum step length
\item[$\Delta_{k}$] SCG comparison parameter
\item[$\Delta t$] LeapFrog time step
\item[$E$] training error
\item[$\varepsilon$] logarithm offset in the cross-entropy
\item[$\mathbf{h}$] hidden activation
\item[$\eta$] SGD learning rate
\item[$\lambda$] SCG scale parameter
\item[$\bar{\lambda}$] SCG scale from the previous iteration
\item[$m$] number of consecutive obtuse LeapFrog angles before $\Delta t$ is halved
\item[$\mu_{k}$] SCG inner product $\mathbf{p}_{k}^{\top}\mathbf{r}_{k}$
\item[$n$] number of training samples
\item[$n_{h}$] number of hidden units
\item[$n_{i}$] number of inputs
\item[$n_{o}$] number of outputs
\item[$N$] number of packed weights
\item[$\mathbf{p}$] SCG search direction
\item[$\mathbf{q}_{k}$] SCG Hessian-vector product
\item[$\mathbf{r}$] steepest-descent vector, $\mathbf{r}=-\nabla E(\mathbf{w})$
\item[$s$] mini-batch size
\item[$\sigma$] SCG finite-difference interval
\item[$\mathbf{u}$] SGD momentum buffer
\item[$\mathbf{v}$] LeapFrog particle velocity
\item[$\mathbf{w}$] packed weight vector
\item[$\mathbf{W}$] weight matrix
\item[$\mathbf{z}$] layer pre-activation
\end{description}

\end{document}
